\section{Veronese 映射}\label{sec:veronese-maps}

本节介绍以意大利数学家 Giuseppe Veronese (1854--1917) 命名的 Veronese 映射. 这是一族射影空间之间的正则映射, 可用固定次数的所有单项式来描述. 在实际应用中, Veronese 映射的作用在于, 它们使我们能够降低射影簇及其间映射的次数. 更确切地说, 如本节将看到的, 可以借助 Veronese 映射证明以下两个令人意外的性质.
\begin{enumerate}[label=\arabic*.]
  \item 在同构意义下, 每个射影簇都可以实现为一族次数至多为二的多项式的消没集 (Proposition~\ref{prop:projective-varieties-defined-by-quadratics}).
  \item 在同构意义下, 射影簇之间的每个多项式映射都是线性映射 (Proposition~\ref{prop:projective-polynomial-maps-linear-up-to-isomorphism}).
\end{enumerate}
现在从定义入手, 开始讨论 Veronese 映射.

\begin{Definition}{Veronese 映射}\leavevmode\label{def:veronese-map}
对任意 $d,n\ge1$, \emph{Veronese 映射} $F_{d,n}:\mathbb{P}^{n}\to\mathbb{P}^{\binom{d+n}{d}-1}$ 是多项式映射
\[
  F_{d,n}=[x_{0}^{d}:x_{0}^{d-1}x_{1}:x_{0}^{d-1}x_{2}:\cdots:x_{n}^{d}],
\]
其中出现在定义中的单项式是变量 $x_{0},\ldots,x_{n}$ 的所有 $d$ 次单项式.
\end{Definition}

Exercise~\ref{ex:count-veronese-monomials-by-stars-and-bars} 概述了一种证明方法, 说明变量 $x_{0},\ldots,x_{n}$ 的 $d$ 次单项式恰有 $\binom{d+n}{d}$ 个 (这解释了作为 $F_{d,n}$ 陪域的射影空间的维数). 此外, Exercise~\ref{ex:veronese-monomials-have-no-common-zero} 要求读者验证, 这些单项式不在 $\mathbb{P}^{n}$ 的任何一点同时消没, 从而解释 $F_{d,n}$ 的定义域为什么是整个 $\mathbb{P}^{n}$.

作为第一个例子, 取 $d=2$ 和 $n=1$, Veronese 映射为
\[
  \begin{aligned}
    F_{2,1}:\mathbb{P}^{1}&\to\mathbb{P}^{2}\\
    [a:b]&\mapsto[a^{2}:ab:b^{2}],
  \end{aligned}
\]
这就是我们在 Examples~\ref{ex:polynomial-map-from-projective-line-to-conic} 和~\ref{ex:projective-line-isomorphic-to-conic} 中考察过的映射 $G$. 在那个特例中, 我们把像 $F_{2,1}(\mathbb{P}^{1})$ 描述为射影簇 $\mathcal{V}(y^{2}-xz)\subseteq\mathbb{P}^{2}$. 下一个目标是把这一描述推广到任意 $d$ 和 $n$ 的 $F_{d,n}(\mathbb{P}^{n})$. 要把任意 Veronese 映射的像描述为射影簇, 首先为 $F_{d,n}$ 陪域的坐标建立方便的记号.

注意, 变量 $x_{0},\ldots,x_{n}$ 的 $d$ 次单项式可以用满足 $d_{0}+\cdots+d_{n}=d$ 的非负整数元组 $D=(d_{0},\ldots,d_{n})$ 来编号. 更确切地说, 与 $D=(d_{0},\ldots,d_{n})$ 对应的单项式为
\[
  x^{D}=x_{0}^{d_{0}}x_{1}^{d_{1}}\cdots x_{n}^{d_{n}}.
\]
据此, 对每个这样的元组 $D$, 将 $\mathbb{P}^{\binom{d+n}{d}-1}$ 的相应坐标记为 $y_{D}$, 从而 $F_{d,n}$ 是 $y_{D}$ 坐标为 $x^{D}$ 的多项式映射.

例如, 在前面考察的 $(d,n)=(2,1)$ 情形下, 将 $\mathbb{P}^{2}$ 的坐标记为 $y_{(2,0)}$, $y_{(1,1)}$ 和 $y_{(0,2)}$. 此时 $F_{2,1}$ 将 $a=[a_{0}:a_{1}]$ 映到 $y_{(2,0)}$ 坐标为 $a_{0}^{2}$, $y_{(1,1)}$ 坐标为 $a_{0}a_{1}$, $y_{(0,2)}$ 坐标为 $a_{1}^{2}$ 的点. 由此可见, $F_{2,1}$ 映到簇
\[
  \mathcal{V}\bigl(y_{(2,0)}y_{(0,2)}-y_{(1,1)}^{2}\bigr)\subseteq\mathbb{P}^{2}.
\]
更一般地, Veronese 映射的像由下面的结果描述.

\begin{Proposition}{Veronese 映射的像}\leavevmode\label{prop:veronese-image-defined-by-quadratics}
对任意 $d,n\ge1$, Veronese 映射 $F_{d,n}:\mathbb{P}^{n}\to\mathbb{P}^{\binom{d+n}{d}-1}$ 的像为
\[
  X_{d,n}
  =\mathcal{V}\bigl(\{y_{D}y_{D'}-y_{E}y_{E'}\mid D+D'=E+E'\}\bigr)
  \subseteq\mathbb{P}^{\binom{d+n}{d}-1}.
\]
\end{Proposition}
\begin{Proof}
$F_{d,n}(\mathbb{P}^{n})\subseteq X_{d,n}$ 来自如下观察: 对任意 $a\in\mathbb{P}^{n}$ 及满足 $D+D'=E+E'$ 的任意元组 $D,D',E,E'$, 有
\[
  a^{D}a^{D'}-a^{E}a^{E'}
  =a_{0}^{d_{0}+d_{0}'}\cdots a_{n}^{d_{n}+d_{n}'}
   -a_{0}^{e_{0}+e_{0}'}\cdots a_{n}^{e_{n}+e_{n}'}=0.
\]
反过来, 为证明 $F_{d,n}(\mathbb{P}^{n})\supseteq X_{d,n}$, 令 $b\in X_{d,n}$. 则对满足 $D+D'=E+E'$ 的所有元组 $D,D',E,E'$, $b$ 的坐标满足 $b_{D}b_{D'}=b_{E}b_{E'}$. 由这些关系可证明 (Exercise~\ref{ex:veronese-coordinate-products-depend-on-sum}) 对满足 $D_{1}+\cdots+D_{k}=E_{1}+\cdots+E_{k}$ 的所有元组 $D_{1},\ldots,D_{k},E_{1},\ldots,E_{k}$, 有
\begin{equation}\label{eq:veronese-equality-of-coordinate-products}
  b_{D_{1}}\cdots b_{D_{k}}=b_{E_{1}}\cdots b_{E_{k}}.
\end{equation}
特别是, 对 $D=(d_{0},\ldots,d_{n})$, 有
\begin{equation}\label{eq:veronese-coordinate-power-from-pure-powers}
  b_{D}^{d}=b_{(d,0,\ldots,0)}^{d_{0}}b_{(0,d,0,\ldots,0)}^{d_{1}}
  \cdots b_{(0,\ldots,0,d)}^{d_{n}}.
\end{equation}
因此坐标 $b_{(d,0,\ldots,0)},b_{(0,d,0,\ldots,0)},\ldots,b_{(0,\ldots,0,d)}$ 中至少一个非零, 否则由 \eqref{eq:veronese-coordinate-power-from-pure-powers} 可得对所有 $D$ 都有 $b_{D}=0$, 这在射影空间中是不可能的.

于是, 不妨设 $b_{(d,0,\ldots,0)}$ 非零, 从而可缩放 $b$ 的齐次坐标, 使 $b_{(d,0,\ldots,0)}=1$. 对每个 $1\le i\le n$, 定义
\[
  a_{i}=b_{(d-1,0,\ldots,0,1,0,\ldots,0)},
\]
其中 $1$ 出现在下标元组的第 $i$ 个坐标. 我们断言
\begin{equation}\label{eq:veronese-preimage-coordinate-reconstruction}
  F_{d,n}(a)=b.
\end{equation}
为证明 \eqref{eq:veronese-preimage-coordinate-reconstruction}, 注意 $F_{d,n}$ 和 $a_{i}$ 的定义意味着, 对任意 $D=(d_{0},\ldots,d_{n})$, $F_{d,n}(a)$ 的 $y_{D}$ 坐标为
\[
  \begin{aligned}
  a_{0}^{d_{0}}a_{1}^{d_{1}}\cdots a_{n}^{d_{n}}
    &=b_{(d,0,\ldots,0)}^{d_{0}}b_{(d-1,1,0,\ldots,0)}^{d_{1}}
      \cdots b_{(d-1,0,\ldots,0,1)}^{d_{n}}\\
    &=b_{(d,0,\ldots,0)}^{d-1}b_{D}=b_{D},
  \end{aligned}
\]
其中第二个等式是 \eqref{eq:veronese-equality-of-coordinate-products} 的应用, 第三个等式来自 $b_{(d,0,\ldots,0)}=1$. 这就证明了断言 \eqref{eq:veronese-preimage-coordinate-reconstruction}, 从而完成了 $X_{d,n}=F_{d,n}(\mathbb{P}^{n})$ 的证明.
\end{Proof}

Proposition~\ref{prop:veronese-image-defined-by-quadratics} 推广了 Example~\ref{ex:projective-line-isomorphic-to-conic} 的一个方面: 它给出了 $F_{d,n}$ 所映入的簇的显式定义方程. 但在 Example~\ref{ex:projective-line-isomorphic-to-conic} 中, 我们实际上还证明了更多: Veronese 映射 $F_{2,1}$ 是到其像上的同构. 这一事实对所有 Veronese 映射也成立.

\begin{Proposition}{Veronese 映射是嵌入}\leavevmode\label{prop:veronese-maps-are-embeddings}
对任意 $d,n\ge1$, Veronese 映射 $F_{d,n}:\mathbb{P}^{n}\to X_{d,n}$ 是同构.
\end{Proposition}
\begin{Proof}
\BookPairedNote{
为证明 $F_{d,n}$ 是到 $X_{d,n}$ 上的同构, 我们构造正则逆映射 $G_{d,n}:X_{d,n}\to\mathbb{P}^{n}$. 给定任意 $b\in X_{d,n}$, 用如下分片描述定义 $G_{d,n}(b)$:
}{
在许多数学语境下,\textquotedbl{}嵌入\textquotedbl{}指到其像上的同构.
}
\[
  G_{d,n}(b)=\left\{\begin{array}{l}
  [b_{(d,0,\ldots,0)}:b_{(d-1,1,\ldots,0)}:\cdots:b_{(d-1,0,\ldots,1)}]\\
  {[b_{(1,d-1,\ldots,0)}:b_{(0,d,\ldots,0)}:\cdots:b_{(0,d-1,\ldots,1)}]}\\
  \multicolumn{1}{c}{\vdots}\\
  {[b_{(1,0,\ldots,d-1)}:b_{(0,1,\ldots,d-1)}:\cdots:b_{(0,0,\ldots,d)}]}
  \end{array}\right.,
\]
各表达式在它给出 $\mathbb{P}^{n}$ 中良定义的点的所有位置都适用. 由 $X_{d,n}$ 的定义方程可见, $G_{d,n}$ 的任意两个表达式的交叉乘积 (见 Exercise~\ref{ex:cross-multiplied-projective-coordinates}) 在 $X_{d,n}$ 上相等, 因此 $G_{d,n}$ 在其有定义的每一点均为良定义. 此外, 如 Proposition~\ref{prop:veronese-image-defined-by-quadratics} 的证明所示, 对每个 $b\in X_{d,n}$, 坐标 $b_{(d,0,\ldots,0)},b_{(0,d,\ldots,0)},\ldots,b_{(0,0,\ldots,d)}$ 中至少一个非零, 从而 $G_{d,n}$ 在所有 $b\in X_{d,n}$ 上都有定义. 因此 $G_{d,n}:X_{d,n}\to\mathbb{P}^{n}$ 是正则映射. 还需核对 $F_{d,n}$ 与 $G_{d,n}$ 是互逆函数, 这留给读者作为习题 (Exercise~\ref{ex:veronese-piecewise-inverse}).
\end{Proof}

簇 $X_{d,n}\subseteq\mathbb{P}^{\binom{n+d}{d}}$ 称为 \emph{Veronese 簇}. Proposition~\ref{prop:veronese-maps-are-embeddings} 表明, 在同构意义下, Veronese 簇 $X_{d,n}$ 就其内在性质而言, 不过是从另一角度看待射影空间 $\mathbb{P}^{n}$. 但研究 $\mathbb{P}^{n}$ 的各种 Veronese 模型, 可以降低射影簇 $X\subseteq\mathbb{P}^{n}$ 的定义多项式的最高次数. 在介绍这一过程一般如何进行之前, 先看一个具体例子.

\begin{Example}{一条三次曲线的 Veronese 像}\leavevmode\label{ex:veronese-image-of-projective-cubic-curve}
令 $X=\mathcal{V}(x_{0}^{3}+x_{1}^{3}+x_{2}^{3})\subseteq\mathbb{P}^{2}$, 考察 Veronese 映射
\[
  F_{3,2}:X\to\mathbb{P}^{\binom{5}{3}-1}=\mathbb{P}^{9}.
\]
回顾 $\mathbb{P}^{9}$ 的坐标 $y_{D}$ 由满足 $d_{0}+d_{1}+d_{2}=3$ 的三元组 $D=(d_{0},d_{1},d_{2})$ 编号. 注意对任意 $a=[a_{0}:a_{1}:a_{2}]\in\mathbb{P}^{2}$, 有
\[
  a\in X\iff a_{0}^{3}+a_{1}^{3}+a_{2}^{3}=0
  \iff F_{3,2}(a)\in\mathcal{V}(y_{(3,0,0)}+y_{(0,3,0)}+y_{(0,0,3)}).
\]
由此得
\[
  F_{3,2}(X)=X_{3,2}\cap\mathcal{V}\bigl(y_{(3,0,0)}+y_{(0,3,0)}+y_{(0,0,3)}\bigr).
\]
此外, 由 Proposition~\ref{prop:veronese-maps-are-embeddings}, $F_{3,2}$ 给出同构 $\mathbb{P}^{2}\cong X_{3,2}$, 从而 $F_{3,2}$ 的限制给出同构
\[
  X\cong X_{3,2}\cap\mathcal{V}\bigl(y_{(3,0,0)}+y_{(0,3,0)}+y_{(0,0,3)}\bigr).
\]
由于 Proposition~\ref{prop:veronese-image-defined-by-quadratics} 表明 $X_{3,2}$ 由二次多项式定义, 这意味着 $X$ 同构于一个可以用次数至多为二的多项式定义的射影簇, 尽管 $X$ 原来的表达式把它描述为一个三次多项式的消没集.
\end{Example}

在前例基础上, 现在描述一般结果.

\begin{Proposition}{射影簇由二次多项式定义}\leavevmode\label{prop:projective-varieties-defined-by-quadratics}
在同构意义下, 每个射影簇都可以写成一族有限个次数至多为二的齐次多项式的消没集.
\end{Proposition}
\begin{Proof}
令 $X\subseteq\mathbb{P}^{n}$ 为射影簇. 首先注意 (Exercise~\ref{ex:choose-projective-defining-polynomials-of-common-degree}), 可以写 $X=\mathcal{V}(f_{1},\ldots,f_{k})$, 其中每个 $f_{i}\in K[x_{0},\ldots,x_{n}]$ 都是同为 $d$ 次的齐次多项式. 换句话说, 可以把每个 $f_{i}$ 写成
\[
  f_{i}=\sum_{D}a_{i,D}x_{0}^{d_{0}}\cdots x_{n}^{d_{n}},
\]
其中求和遍历所有元素之和为 $d$ 的非负整数元组 $D=(d_{0},\ldots,d_{n})$, 且 $a_{i,D}\in K$. 考察 Veronese 映射
\[
  F_{d,n}:\mathbb{P}^{n}\to\mathbb{P}^{\binom{n+d}{d}-1},
\]
用陪域的变量定义线性多项式 $\ell_{1},\ldots,\ell_{k}$:
\[
  \ell_{i}=\sum_{D}a_{i,D}y_{D}.
\]
如 Example~\ref{ex:veronese-image-of-projective-cubic-curve} 所示, 对任意 $a\in\mathbb{P}^{n}$, 有
\[
  a\in X\iff F_{d,n}(a)\in X_{d,n}\cap\mathcal{V}(\ell_{1},\ldots,\ell_{k}),
\]
从而 $F_{d,n}(X)=X_{d,n}\cap\mathcal{V}(\ell_{1},\ldots,\ell_{k})$. 由于 $F_{d,n}:\mathbb{P}^{n}\to\mathbb{P}^{\binom{n+d}{d}-1}$ 是正则映射, 它可限制为正则映射
\[
  F_{d,n}:X\to X_{d,n}\cap\mathcal{V}(\ell_{1},\ldots,\ell_{k}).
\]
又因为 $F_{d,n}$ 是到 $X_{d,n}$ 上的同构 (Proposition~\ref{prop:veronese-maps-are-embeddings}), 所以 $F_{d,n}^{-1}:X_{d,n}\to\mathbb{P}^{n}$ 可限制为正则逆映射
\[
  F_{d,n}^{-1}:X_{d,n}\cap\mathcal{V}(\ell_{1},\ldots,\ell_{k})\to X.
\]
因此, $X\cong X_{d,n}\cap\mathcal{V}(\ell_{1},\ldots,\ell_{k})$. 由 $X_{d,n}$ 可以用二次多项式定义 (Proposition~\ref{prop:veronese-image-defined-by-quadratics}), 而各 $\ell_{i}$ 都是线性的, 结论随即成立.
\end{Proof}

值得注意的是, Proposition~\ref{prop:projective-varieties-defined-by-quadratics} 虽然降低了射影簇的定义多项式的最高次数, 却通常会使定义多项式的数量大幅增加. 例如, 在 Example~\ref{ex:veronese-image-of-projective-cubic-curve} 中, 可以利用维数理论基本定理的一个射影版本 (将在本书最后一章中精确说明), 证明描述这条三次曲线的 Veronese 像至少需要 8 个线性和二次多项式, 虽然三次曲线本身只需一个定义多项式.

作为 Veronese 映射的最后一个应用, 下一个结果说明, 在对定义域作同构的意义下, 所有多项式映射都是线性映射.

\begin{Proposition}{在同构意义下, 多项式映射是线性的}\leavevmode\label{prop:projective-polynomial-maps-linear-up-to-isomorphism}
若 $F:X\to Y$ 是射影簇之间的多项式映射, 则存在同构 $G:X\to Z$ 及线性映射 $L:Z\to Y$, 使 $F=L\circ G$.
\end{Proposition}
\begin{Proof}
Exercise~\ref{ex:linearize-projective-polynomial-map-by-veronese}.
\end{Proof}

Proposition~\ref{prop:projective-polynomial-maps-linear-up-to-isomorphism} 的一个令人意外的推论, 是下述关于多项式映射的有限对一结果. 这再次表明多项式映射的限制有多强, 也进一步说明我们一直研究的较灵活的正则映射概念为何合理: 正则映射只在局部是多项式的.

\begin{Corollary}{多项式映射是有限对一的}\leavevmode\label{cor:projective-polynomial-maps-are-finite-to-one}
若 $F:X\to Y$ 是射影簇之间的多项式映射, 则对任意 $b\in Y$, 使 $F(a)=b$ 的 $a\in X$ 只有有限个.
\end{Corollary}
\begin{Proof}
沿用 Proposition~\ref{prop:projective-polynomial-maps-linear-up-to-isomorphism} 陈述中的记号, 任意多项式映射 $F$ 都可写成复合 $L\circ G$, 其中 $G$ 是同构, $L$ 是线性的. 由同构是一对一的, 而射影簇之间的线性映射是有限对一的 (Exercise~\ref{ex:normal-form-and-finite-fibres-of-projective-linear-map}), 即得结论.
\end{Proof}

Exercise~\ref{ex:regular-projection-from-projective-quadric} 给出了一个不是有限对一的正则映射的例子. 于是, 由 Corollary~\ref{cor:projective-polynomial-maps-are-finite-to-one}, 这个正则映射不是多项式映射.

\BookExerciseHeading
\begin{BookExercises}
  \item\label{ex:count-veronese-monomials-by-stars-and-bars} 设 $d,n\ge1$ 为整数, 令 $SB_{d,n}$ 为由 $d$ 个\textquotedbl{}星号\textquotedbl{}和 $n$ 个\textquotedbl{}竖线\textquotedbl{}组成的序列的集合. 例如, 当 $(d,n)=(4,3)$ 时, $SB_{d,n}$ 的几个元素为
  \[
    (\star\mid\star\star\mid\mid\star),\qquad
    (\mid\mid\mid\star\star\star\star),\qquad
    (\star\mid\star\mid\star\mid\star).
  \]
  \begin{enumerate}
    \item 解释为什么 $SB_{d,n}$ 有 $\binom{d+n}{d}$ 个元素.
    \item 描述 $SB_{d,n}$ 和变量 $x_{0},\ldots,x_{n}$ 的 $d$ 次单项式的集合之间的双射.
    \item 由 (a) 和 (b) 推出变量 $x_{0},\ldots,x_{n}$ 的 $d$ 次单项式恰有 $\binom{d+n}{d}$ 个.
  \end{enumerate}
  \item\label{ex:veronese-monomials-have-no-common-zero} 证明第 $d$ 个 Veronese 映射 $F_{d,n}$ 在 $\mathbb{P}^{n}$ 的每一点都正则.
  \item\label{ex:veronese-coordinate-products-depend-on-sum} 令 $X_{d,n}\subseteq\mathbb{P}^{\binom{d+n}{d}-1}$ 为 Veronese 簇, 考察一点 $b\in X_{d,n}$. 本题证明坐标的乘积
  \[
    b_{D_{1}}\cdots b_{D_{k}}\in K
  \]
  只依赖于 $D_{1}+\cdots+D_{k}\in\mathbb{N}^{n+1}$. 固定一个满足 $e_{0}+\cdots+e_{n}=dk$ 的元组 $(e_{0},\ldots,e_{n})\in\mathbb{N}^{n+1}$. 令 $\mathcal{S}$ 表示长度为 $dk$ 的序列的集合, 其中有 $e_{0}$ 个元素等于 $0$, 有 $e_{1}$ 个元素等于 $1$, 依此类推. 对每个 $\sigma\in\mathcal{S}$, 令
  \[
    b_{\sigma}=b_{D_{1}}\cdots b_{D_{k}},
  \]
  以定义 $b_{\sigma}\in K$, 其中 $D_{j}$ 的第 $i$ 项为 $\sigma$ 的第 $j$ 个长度为 $d$ 的子序列中 $i$ 的个数.
  \begin{enumerate}
    \item 为理解记号, 考察元组 $(e_{0},e_{1},e_{2})=(3,2,4)$, 这里取 $n=2$ 和 $d=k=3$. 写出三个序列 $\sigma\in\mathcal{S}$ 的例子及相应的值 $b_{\sigma}\in K$.
    \item 证明将 $\sigma$ 送到 $b_{\sigma}$ 的函数 $\phi:\mathcal{S}\to K$ 满射到
    \[
      \{b_{D_{1}}\cdots b_{D_{k}}\mid D_{1}+\cdots+D_{k}=(e_{1},\ldots,e_{n})\}.
    \]
    \item 设 $\sigma$ 与 $\sigma'$ 是 $\mathcal{S}$ 中仅相差两个元素对换的序列. 用 $X_{d,n}$ 的定义方程证明
    \[
      \phi(\sigma)=\phi(\sigma').
    \]
    \item 用对换生成所有排列这一事实, 推出 $b_{D_{1}}\cdots b_{D_{k}}$ 只依赖于 $D_{1}+\cdots+D_{k}$.
  \end{enumerate}
  \item\label{ex:veronese-piecewise-inverse} 沿用 Proposition~\ref{prop:veronese-maps-are-embeddings} 证明中的记号, 验证 $F_{d,n}$ 和 $G_{d,n}$ 是互逆函数.
  \item\label{ex:choose-projective-defining-polynomials-of-common-degree}
  \begin{enumerate}
    \item 令 $f\in K[x_{0},\ldots,x_{n}]$ 且 $d\in\mathbb{N}$. 证明
    \[
      \mathcal{V}_{\mathbb{P}}(f)=\mathcal{V}_{\mathbb{P}}(x_{0}^{d}f,\ldots,x_{n}^{d}f).
    \]
    \item 证明每个射影簇都可以用一组有限个同次数的齐次多项式定义.
  \end{enumerate}
  \item\label{ex:linearize-projective-polynomial-map-by-veronese} 证明 Proposition~\ref{prop:projective-polynomial-maps-linear-up-to-isomorphism}. (提示: 取 $G$ 为一个 Veronese 映射.)
\end{BookExercises}
